\paragraph{A school admissions system with reported preferences and initial offers.}
Chile's Ministry of Education releases the regular-round lists families submitted, offers before waitlists, program vacancies, priority indicators, and school lottery orders \citep{saedata,correaschool}. We trained on the 2019 application cycle and tested on every applicant in 2020 and 2021. For each applicant we compare shortlists of five reported courses, keeping an entitled current-school option as an additional fallback. The model predicts the 2019 official offer using capacities, aggregate submitted demand, priority tiers and prior program offer frequencies. At test time it uses the reported choice set but cannot see its order or the realized lottery; an additional model receives the preference positions. The replay uses the actual reported order, lottery and coarse priorities. It pools different seat types and freezes family priorities, reproducing the exact official course or unassigned status for 95.24\% and 95.52\% of applicants (Appendix~\ref{app:sae}). Only coverage of the historical offer is observed-label evidence. The number who would be placed after a menu change is counterfactual.

The original contrast changes sign. Against uniform menus, the rank-blind model improves official-offer coverage by 1.80 and 1.52 percentage points and \emph{increases} placement by 0.28 and 0.24 points. A different comparison exposes the prediction cost. Adding preference positions increases official-offer coverage by a further 0.32 and 0.33 points, while costing 525 and 519 placements. At a single discretionary slot the same comparison gains 9.64 and 9.95 points of coverage and loses 3.73 and 3.39 points of placement. Among applicants whose own menus never change at quota five, the rank-blind policy reduces placement by 0.13 and 0.07 points relative to uniform, even though it improves placement overall.
